%%==========================================================================%%
%% datarecords.tex -- "Data Records".                                         %%
%% The repository, files, formats, folder structure, and definitions of any   %%
%% non-obvious fields. NO summary statistics here. Host is [TBD], so both      %%
%% the DICOM and the file-based layouts are described.                        %%
%% American spelling.                                                         %%
%%==========================================================================%%

\section*{Data Records}\label{sec:datarecords}

The dataset is deposited in [TBD: repository] under accession [TBD] and Digital Object Identifier [TBD: https://doi.org/XXXX].
The imaging, the segmentation labels, and the accompanying metadata are organized as described below.
The host repository determines the released formats, and both candidate layouts are described because the host is not yet fixed.

\subsection*{File organization}\label{subsec:layout}
On a DICOM host, the imaging follows the standard DICOM hierarchy.
Each subject contains one or more imaging studies, and each study contains the MRI series acquired in that session.
Studies and series are identified by their DICOM unique identifiers, which are remapped during de-identification, and a per-subject date offset preserves the relative timing between a subject's examinations.
Segmentation labels are stored as DICOM-SEG objects that reference the corresponding image series.

On a file-based host, the imaging is provided as NIfTI volumes and the segmentation labels as NIfTI or NRRD label volumes, organized by subject in a directory hierarchy that mirrors the DICOM grouping.
The label volumes are spatially aligned to the corresponding imaging.
A representative directory layout for the file-based release is shown below [TBD: finalize to match the chosen host].

\begin{verbatim}
dataset_root/
  sub-[ID]/
    ses-[ID]/
      anat/
        sub-[ID]_ses-[ID]_[sequence].nii.gz   # whole-body MRI series
      seg/
        sub-[ID]_ses-[ID]_anatomy.nii.gz       # multi-structure label volume
  metadata/
    examinations.csv                            # per-examination metadata
    label_index.csv                             # label-index map (data dictionary)
    data_dictionary.csv                         # column definitions
\end{verbatim}

\subsection*{Imaging data}\label{subsec:imaging}
The MRI image data are provided in [TBD: DICOM or NIfTI], with the acquisition parameters retained in the DICOM headers and mirrored in the per-examination metadata.

\subsection*{Segmentation labels and label-index map}\label{subsec:labels}
The released segmentation labels cover 13 anatomical structures.
Seven structures are bilateral and stored as separate left and right instances, while six are single structures, giving 20 labels in total.
The labels are organized in the deposited directory tree as [TBD: folder grouping to match the deposited directory tree].
The mapping between label index, structure, and laterality is the central data dictionary for the dataset and is provided as a machine-readable file with the release.
Table~\ref{tab:labels} summarizes the structures and their organization.

\begin{table}[h]
\caption{Structures covered by the released anatomy segmentations. Bilateral structures are stored as left and right instances. The complete label-index map is provided as a machine-readable file with the release [TBD: confirm the exact index values and folder names against the deposited files].}\label{tab:labels}
\begin{tabular}{@{}lll@{}}
\toprule
Group & Structure & Laterality \\
\midrule
Skeleton & Femur          & Left, Right \\
Skeleton & Tibia          & Left, Right \\
Skeleton & Fibula         & Left, Right \\
Skeleton & Humerus        & Left, Right \\
Skeleton & Hip            & Left, Right \\
Skeleton & Spine          & Single \\
Skeleton & Sacrum         & Single \\
Organ    & Lung           & Left, Right \\
Organ    & Kidney         & Left, Right \\
Organ    & Liver          & Single \\
Organ    & Spleen         & Single \\
Organ    & Heart          & Single \\
Organ    & Urinary bladder & Single \\
\botrule
\end{tabular}
\end{table}

\subsection*{Metadata}\label{subsec:metadata}
A per-examination metadata table is provided as an open plain-text file (CSV) with an accompanying data dictionary, so the dataset is usable with standard tooling.
The table captures the indexing variables useful for cohort selection and modeling, grouped as in Table~\ref{tab:vars}.
The annotation protocol is recorded per labeled examination, distinguishing the correction-based labels from the from-scratch labels (Section~\ref{sec:methods}).

\begin{table}[h]
\caption{Indexing variables provided with the dataset [TBD: finalize against the deposited metadata files].}\label{tab:vars}
\begin{tabular}{@{}ll@{}}
\toprule
Group & Example fields \\
\midrule
Subject and demographics & subject\_id, age\_at\_scan, sex \\
Clinical context         & diagnosis, lymphoma\_subtype, protocol\_provenance \\
Acquisition              & scanner, field\_strength, sequence, voxel\_spacing, TR, TE \\
Labels and annotation    & label\_id, structure, laterality, annotation\_protocol \\
File and quality control & file\_reference, checksum, qc\_status, license \\
\botrule
\end{tabular}
\end{table}
